\section{Wyniki pomiarów}
\label{sec:wyniki}
W tabeli~\ref{tab:wyniki} zestawiono wszystkie dostępne próbki
zarejestrowane miernikiem. Każda próbka trwała $10\,\mathrm{s}$.
Wartości odnoszą się do profilu A i liniowego całkowania Leq.
\begin{table}[htbp]
 \centering
 \caption{Równoważne poziomy dźwięku A w próbkach $10\,\mathrm{s}$.}
 \label{tab:wyniki}
 \begin{tabular}{@{}lrlr@{}}
\toprule
Plik: stacja & $L_{Aeq,10s}$ [dB] & Plik: tło & $L_{Aeq,10s}$ [dB] \\
\midrule
P04 & 52,54 & P10 & 48,05 \\
P05 & 52,27 & P11 & 47,72 \\
P06 & 51,96 & P12 & 49,77 \\
P07 & 53,31 & P13 & 47,08 \\
P08 & 53,42 & P14 & 47,20 \\
\bottomrule
\end{tabular}

\end{table}


\begin{table}[htbp]
 \centering
 \caption{Opracowanie serii i sprawdzenie wymaganej liczby próbek.}
 \label{tab:opracowanie}
 \begin{tabular}{@{}lrr@{}}
\toprule
Wielkość & Przy stacji \\
\midrule
Średnia energetyczna [dB] & 52,74 \\
Rozstęp $R$ [dB] & 1,46  \\
Liczba wykonanych próbek & 5  \\
Wymagana liczba próbek & 4  \\
Kryterium liczby próbek & spełnione  \\
\bottomrule
\end{tabular}

\end{table}
Seria zasadnicza spełnia kryterium liczby próbek. 

Różnica średnich poziomów wyniosła $\Delta L=4,66\,\mathrm{dB}$.
Po podstawieniu do wzoru~\eqref{eq:korekcja-tla} otrzymano
\begin{equation}
 L_{Aek}=50,92\,\mathrm{dB}.
\end{equation}

Ze względu na brak sprawdzenia toru
kalibratorem oraz brak pełnej dokumentacji warunków wynik ten
należy traktować jako oszacowanie dla warunków ćwiczenia.
Nie stanowi on podstawy do stwierdzenia dotrzymania lub przekroczenia
poziomu dopuszczalnego na terenie chronionym przed hałasem.
